\section{Conclusiones}
\label{sec:conclusiones}
% Síntesis, respuesta a los objetivos, recomendaciones, líneas futuras. Sin resultados nuevos.

Los tres experimentos responden al objetivo general con un mismo criterio. Las propiedades
algebraicas de la suma dejan de valer en la computadora cuando la magnitud del sumando y la del
acumulador se separan lo suficiente como para que el término chico caiga por debajo de la resolución
local del formato. El resto son variantes de esa situación: el 1 que desaparece frente a $2^{53}$,
los términos tardíos de una serie decreciente que se pierden contra un acumulador cercano a 1, y las
cifras de una raíz que la resta deja al descubierto. La condición depende del espaciado de los
números representables, no del lenguaje ni de la implementación.

Sobre el primer objetivo específico, el umbral de la pérdida quedó determinado y verificado. Con
punto flotante de doble precisión, el término $1 + b^{k} - b^{k}$ deja de aportar pasado el mayor
$k$ con $|b|^{k} < 2^{53}$, y ese valor predicho coincidió con la meseta medida en seis de las ocho
bases ensayadas, con un término de diferencia en las dos restantes.
La agrupación de los paréntesis solo cambia el resultado cuando la potencia cae exactamente en
$2^{53}$, el borde donde el espaciado pasa de 1 a 2; entre las bases ensayadas eso ocurre
únicamente con $b = \pm 2$. El
entero de precisión arbitraria de Python conserva la identidad para todo $N$, al costo de un tiempo
que crece de manera cuadrática. El entero de ancho fijo también la conserva, pero por una razón
distinta: la reducción módulo $2^{64}$ preserva la estructura de anillo y la resta deshace el
desborde. Esa coincidencia no vuelve seguro al entero de ancho fijo, porque depende de que el valor
desbordado se cancele dentro de la misma expresión.

Sobre el segundo, el orden de acumulación resultó ser un factor de primer orden y no un detalle de
implementación. En la misma serie, los cuatro algoritmos se separaron en dos grupos que difieren en
más de dos órdenes de magnitud de error en el rango grande, y las 50 permutaciones aleatorias de un
mismo cálculo produjeron
48 resultados distintos. La conclusión práctica es más específica de lo previsto:
acumular de menor a mayor módulo alcanzó para igualar a la suma compensada de Kahan en esta serie,
de modo que cuando el problema es el orden, ordenar ya lo resuelve. Kahan conserva la ventaja
cuando el orden de las magnitudes no se conoce sin calcularlas.

Sobre el tercero, la comparación entre las dos sucesiones delimita cuándo la cancelación entre
cantidades próximas es un problema. En $b_N$ la representación telescópica resta términos cada vez
más parecidos y aun así no se degrada, porque lo que pierde son cifras bajas de un sumando que ya
era pequeño frente al acumulado. En $c_N$ la misma transformación destruye el resultado, porque ahí
el término cancelado es el resultado mismo. El fenómeno se entiende mejor si se separa el lema de
Sterbenz del nombre que lleva: la resta suele ser exacta y no crea error, solo expone el que los
operandos traían de un redondeo anterior.

Sobre el cuarto, las cotas teóricas describieron los datos con distinta fidelidad. La referencia
estadística $\sqrt{N}\,u$ contuvo a las curvas de error en todo el rango ensayado. La cota de peor
caso $Nu$ las sobreestimó por más de tres órdenes de magnitud, lo que confirma que suponer redondeos
alineados describe un caso posible y no el típico. La cota $2k^{2}u$ de la cancelación resultó la
más precisa de las tres: predijo los umbrales medidos
dentro de un factor 1.33 y acertó el valor exacto del $k$ en que el término colapsa. Contar
operaciones, en cambio, acota el error sin predecirlo caso por caso, como muestran las dos formas
cerradas, cuyas cotas difieren en un factor 2 y cuyos resultados casi siempre coinciden.

Sobre el quinto, cada figura y cada tabla del informe se genera con una sola corrida del código del
repositorio, y ningún número del texto proviene de la consola.

De todo esto se desprenden cuatro recomendaciones para quien programe una sumatoria larga. Sumar de
menor a mayor módulo cuando el orden de las magnitudes se conoce de antemano, y recurrir a una suma
compensada cuando no. Preferir, entre expresiones algebraicamente equivalentes, la que evita restar
cantidades próximas, que en el caso estudiado se obtiene racionalizando por el conjugado. Fijar y
documentar la semilla si el orden de acumulación depende de un sorteo, porque de lo contrario el
cálculo no es reproducible. Y verificar el resultado, ya que ninguno de los fenómenos estudiados
interrumpe el programa: el desborde devuelve $\infty$ y \texttt{NaN} en silencio, y el entero de
ancho fijo puede entregar el valor correcto y ocultar que las operaciones intermedias desbordaron.

Quedan abiertas varias líneas. Repetir las mediciones en \texttt{binary32} permitiría comprobar si
los umbrales se desplazan como predicen las fórmulas. Usar precisión extendida como patrón separaría
el error del algoritmo del error de la referencia, que en este trabajo comparten el mismo piso.
Extender el estudio a series con términos de signos alternados pondría a prueba la recomendación
sobre el orden, que fue establecida sobre series de términos positivos y podría invertirse. Falta
también comparar la acumulación secuencial contra la suma por pares y contra variantes más
elaboradas de compensación, que Higham (2002) analiza y que este trabajo no ensayó. Queda por medir
qué ocurre cuando el compilador o la paralelización reasocian las operaciones por su cuenta, un
escenario donde el programador ya no controla el agrupamiento.
